% Source: https://artofproblemsolving.com/community/c3h2299958_angle_wanted_36o_twice__2014_estonia_open_junior_22
% Stable identity: aops-topic-2299958-q01
Let the two angles in the ratio $1:2$ be $\alpha$ and $2\alpha$, and let the third angle be $\beta=180^\circ-3\alpha$. A scalene triangle has three different sides, hence three different angles. One of the three angles is $36^\circ$:
\begin{itemize}
\item $\alpha=36^\circ$: the angles are $36^\circ$, $72^\circ$, $72^\circ$. Two angles are equal, so the triangle is isosceles; rejected.
\item $2\alpha=36^\circ$: $\alpha=18^\circ$ and $\beta=126^\circ$, giving $18^\circ$, $36^\circ$, $126^\circ$, all different.
\item $\beta=36^\circ$: $3\alpha=144^\circ$, so $\alpha=48^\circ$, giving $48^\circ$, $96^\circ$, $36^\circ$, all different.
\end{itemize}

\begin{asy}
usepackage("amsmath");
size(10cm);
defaultpen(fontsize(10pt));
void am(pair V, pair P, pair Q, string s, real r){ pair w=unit(unit(P-V)+unit(Q-V)); label(s,V+r*w,fontsize(8pt)); }
pair A=(0,0), B=(7,0);
pair C=extension(A,A+dir(18),B,B+dir(180-36));
draw(A--B--C--cycle,linewidth(0.9));
am(A,B,C,"$18^\circ$",2.1); am(B,A,C,"$36^\circ$",1.3); am(C,A,B,"$126^\circ$",0.75);
pair A2=(9,0), B2=(16,0);
pair C2=extension(A2,A2+dir(48),B2,B2+dir(180-36));
draw(A2--B2--C2--cycle,linewidth(0.9));
am(A2,B2,C2,"$48^\circ$",1.1); am(B2,A2,C2,"$36^\circ$",1.3); am(C2,A2,B2,"$96^\circ$",0.9);
label("$36^\circ=2\cdot18^\circ$",(3.5,-0.8),fontsize(8pt)); label("$96^\circ=2\cdot48^\circ$",(12.5,-0.8),fontsize(8pt));
\end{asy}

So the angles are either $\boxed{18^\circ,\ 36^\circ,\ 126^\circ}$ or $\boxed{36^\circ,\ 48^\circ,\ 96^\circ}$. (Both triangles exist, since the angles are positive and add up to $180^\circ$, and they are scalene because their angles are different.)
